% 基于 EmpyreanHYR 双栏中文论文模板的版式思想重构。
% 上游许可与版本见 licenses/upstream_template.md 及
% licenses/upstream_template_LICENSE.txt。

% ---------- 字体：优先课程目标字体，缺失时自动后备 ----------
\usepackage{fontspec}
\usepackage{xeCJK}

\IfFontExistsTF{Times New Roman}{
  \setmainfont{Times New Roman}
}{
  \setmainfont{TeX Gyre Termes}
}

\IfFontExistsTF{SimSun}{
  \setCJKmainfont{SimSun}[BoldFont=SimHei]
}{
  \setCJKmainfont{FandolSong}[BoldFont=FandolSong-Bold]
}
\IfFontExistsTF{SimHei}{
  % SimHei 通常只有常规字形；启用受控伪粗体以避免标题层级调用
  % \bfseries 时发生字形替代警告。
  \setCJKsansfont{SimHei}[AutoFakeBold=2]
  \setCJKfamilyfont{zhhei}{SimHei}[AutoFakeBold=2]
}{
  \setCJKsansfont{FandolHei}[BoldFont=FandolHei-Bold]
  \setCJKfamilyfont{zhhei}{FandolHei}[BoldFont=FandolHei-Bold]
}
\IfFontExistsTF{FandolFang}{
  \setCJKfamilyfont{zhfangsong}{FandolFang}[AutoFakeBold=2]
}{
  \setCJKfamilyfont{zhfangsong}{FandolSong}[BoldFont=FandolSong-Bold]
}
\setCJKmonofont{FandolFang}[AutoFakeBold=2]
\newcommand{\fangsongfont}{\CJKfamily{zhfangsong}}
\newcommand{\heiti}{\CJKfamily{zhhei}}

% ---------- 页面与双栏 ----------
\usepackage[left=2cm,right=2cm,top=2cm,bottom=2cm,
            headheight=14pt,headsep=0.45cm,footskip=0.8cm]{geometry}
\setlength{\columnsep}{0.8cm}
\setlength{\columnseprule}{0pt}
\setlength{\parindent}{2em}
\setlength{\parskip}{0pt}
\linespread{1.08}
\usepackage{indentfirst}
\usepackage{setspace}
\usepackage{ragged2e}
\usepackage{microtype}
\XeTeXlinebreaklocale "zh"
\XeTeXlinebreakskip = 0pt plus 1pt
\widowpenalty=10000
\clubpenalty=10000
\displaywidowpenalty=10000
\raggedbottom

% ---------- 标题层级 ----------
\ctexset{
  section={
    format=\heiti\bfseries\zihao{-4},
    number=\arabic{section},
    aftername=\hspace{0.55em},
    beforeskip=0.85em,
    afterskip=0.35em,
    indent=0pt
  },
  subsection={
    format=\heiti\bfseries\zihao{5},
    number=\thesection.\arabic{subsection},
    aftername=\hspace{0.5em},
    beforeskip=0.65em,
    afterskip=0.25em,
    indent=0pt
  },
  subsubsection={
    format=\heiti\zihao{5},
    number=\thesubsection.\arabic{subsubsection},
    aftername=\hspace{0.45em},
    beforeskip=0.5em,
    afterskip=0.2em,
    indent=0pt
  }
}

% ---------- 数学 ----------
\usepackage{amsmath,amssymb,mathtools,bm}
\numberwithin{equation}{section}

% ---------- 图表 ----------
\usepackage{graphicx}
\usepackage{booktabs,tabularx,array,threeparttable,multirow,makecell}
\usepackage{caption}
\usepackage{subcaption}
\usepackage{tikz}
\usetikzlibrary{arrows.meta,positioning,calc,fit}
\usepackage{siunitx}
\sisetup{detect-all}
\usepackage{needspace}
\usepackage{enumitem}

\renewcommand{\figurename}{图}
\renewcommand{\tablename}{表}
\captionsetup{
  font={small},
  labelfont=bf,
  labelsep=quad,
  justification=centering,
  singlelinecheck=true,
  skip=4pt
}

% ---------- 页眉页脚 ----------
\usepackage{fancyhdr}
\pagestyle{fancy}
\fancyhf{}
\fancyhead[L]{\zihao{-5}\fangsongfont 无线边缘网络中联邦学习通信效率优化方法研究综述}
\fancyhead[R]{\rmfamily\small\thepage}
\renewcommand{\headrulewidth}{0.4pt}
\renewcommand{\footrulewidth}{0pt}

% ---------- GB/T 7714—2015 顺序编码参考文献 ----------
\usepackage[
  backend=biber,
  style=gb7714-2015,
  sorting=none,
  giveninits=true,
  gbpub=false,
  gbnamefmt=lowercase
]{biblatex}
\addbibresource{references.bib}
\setlength{\bibhang}{2em}
\setlength{\bibitemsep}{0.05em}
\renewcommand*{\bibfont}{\zihao{5}\selectfont\setstretch{1.0}}
\AtBeginBibliography{\RaggedRight\setlength{\emergencystretch}{2em}}

% ---------- 超链接与中文交叉引用 ----------
\usepackage[hidelinks,unicode,bookmarks=true,bookmarksnumbered=true]{hyperref}
\usepackage[nameinlink,noabbrev]{cleveref}
\crefname{figure}{图}{图}
\Crefname{figure}{图}{图}
\crefname{table}{表}{表}
\Crefname{table}{表}{表}
\crefname{equation}{式}{式}
\Crefname{equation}{式}{式}
\crefname{section}{第}{第}
\Crefname{section}{第}{第}
\crefname{subsection}{第}{第}
\Crefname{subsection}{第}{第}

% ---------- 常用列类型 ----------
\newcolumntype{Y}{>{\centering\arraybackslash}X}
\newcolumntype{L}{>{\raggedright\arraybackslash}X}
